\lysh{18244_SOLUTION}{1}
\textbf{ΛΥΣΗ}

α) Η $\varepsilon_1$ διέρχεται από την αρχή των αξόνων και από το σημείο $\text{A}(1,\sqrt{3})$ ενώ η $\varepsilon_2$ διέρχεται από την αρχή των αξόνων και από το σημείο $\text{B}(1,1)$, όπως φαίνεται στο παρακάτω σχήμα.

β) Η $\varepsilon_1$ έχει συντελεστή διεύθυνσης $\sqrt{3}$ οπότε σχηματίζει με τον $\text{xx}'$ γωνία $60^\circ$, ενώ η $\varepsilon_2$ έχει συντελεστή διεύθυνσης $1$ οπότε σχηματίζει με τον $\text{xx}'$ γωνία $45^\circ$.

γ) Όπως φαίνεται και στο παρακάτω σχήμα, η οξεία γωνία που σχηματίζουν οι ευθείες $\varepsilon_1,\varepsilon_2$ είναι η διαφορά των γωνιών που σχηματίζει η κάθε μία από τις $\varepsilon_1,\varepsilon_2$ με τον $\text{xx}'$, δηλαδή $60^\circ-45^\circ=15^\circ$.

δ) Το $\vec{\delta_1} = (1,\sqrt{3})$ είναι παράλληλο στην ευθεία $\varepsilon_1$ και το $\vec{\delta_2} = (1,1)$ είναι παράλληλο στην ευθεία $\varepsilon_2$.
Είναι
\[
\left|\vec{\delta_1}\right| = \sqrt{1^2+(\sqrt{3})^2} = \sqrt{1+3} = \sqrt{4} = 2, \quad \left|\vec{\delta_2}\right| = \sqrt{1^2+1^2} = \sqrt{2}
\]
και
\[
\vec{\delta_1} \cdot \vec{\delta_2} = (1,\sqrt{3}) \cdot (1,1) = 1+\sqrt{3} > 0
\]
οπότε
\[
\sigma\upsilon\nu 15^\circ = \sigma\upsilon\nu \left(\widehat{\vec{\delta_1},\vec{\delta_2}}\right) = \frac{\vec{\delta_1} \cdot \vec{\delta_2}}{|\vec{\delta_1}| \cdot |\vec{\delta_2}|} = \frac{1+\sqrt{3}}{2\sqrt{2}}.
\]

% TODO: TikZ σχήμα
\endlysh